\chapter{Grover's Search Algorithm}

% The braket package ships \ketbra commented out; define it if the preamble has not.
\providecommand{\ketbra}[2]{\ket{#1}\!\bra{#2}}

Grover's algorithm solves the \emph{unstructured search} problem with a
quadratic speedup over any classical method. Given oracle access to a Boolean
function that flags a unique marked item among $N=2^{n}$ candidates, a classical
computer must, in the worst case, probe essentially all of them; Grover's
algorithm finds the marked item using only $O(\sqrt N)$ oracle queries. The
entire mechanism reduces to a strikingly simple geometric picture: the search
takes place inside a two-dimensional plane, and each iteration is a rotation
toward the solution. This chapter builds that picture from the ground up, proves
the key facts, computes the optimal number of iterations, and closes with the
generalization to amplitude amplification and the matching $\Omega(\sqrt N)$
lower bound that shows the algorithm is optimal.

\section{The Unstructured Search Problem}

We begin by fixing the problem and the model of computation. The word
\emph{unstructured} is essential: there is no order, index, or exploitable
regularity among the items, so the trick behind binary search is unavailable.

\begin{definition}[Unstructured search]
Let $N=2^{n}$ and let $f\colon\{0,1\}^{n}\to\{0,1\}$ be a Boolean function to
which we have only \emph{oracle} access. We are promised that there is a unique
marked item $x_0\in\{0,1\}^{n}$ with
\[
  f(x)=\begin{cases}1,& x=x_0,\\ 0,& x\neq x_0.\end{cases}
\]
The oracle is the unitary $\Uf$ acting on $n+1$ qubits by
\[
  \Uf\ket{x}\ket{y}=\ket{x}\ket{y\oplus f(x)},
  \qquad x\in\{0,1\}^{n},\ y\in\{0,1\}.
\]
The goal is to determine $x_0$ using as few queries to $\Uf$ as possible.
\end{definition}

Each application of $\Uf$ counts as one \emph{query}, and we measure the cost of
an algorithm by its query complexity. Because the function has no structure, the
oracle is the only way to learn anything about $x_0$: the string $x_0$ cannot be
read off directly.

\begin{proposition}[Classical lower bound]
Any classical (possibly randomized) algorithm that identifies $x_0$ with constant
success probability must make $\Omega(N)$ queries to $f$ in the worst case.
\end{proposition}

\begin{proof}[Sketch]
A classical algorithm learns nothing about $x_0$ from a query $f(x)$ unless
$x=x_0$. An adversary may answer $f(x)=0$ for every queried $x$ until only one
unqueried candidate remains, which is consistent with $x_0$ being that last item.
Thus to be certain, an algorithm must probe $N-1$ items in the worst case. A
randomized algorithm querying a uniformly random subset of size $k$ finds $x_0$
with probability $k/N$, so achieving constant success probability requires
$k=\Omega(N)$ queries.
\end{proof}

This abstract task captures many concrete situations: searching an unsorted
database, deciding satisfiability of a constraint problem (treat any satisfying
assignment as a marked item and let the oracle test assignments), or inverting a
function and finding collisions. In every case the common structure is the same:
\emph{package ``is this candidate a solution?'' into an oracle}, and the classical
exhaustive cost is $O(N)$. Grover's algorithm attacks all of them at once.

\subsection{The phase oracle}

It is convenient to convert the bit-flip oracle $\Uf$ into a \emph{phase} oracle.
Prepare the answer qubit in the state
$\ket{-}=\tfrac{1}{\sqrt2}\bigl(\ket{0}-\ket{1}\bigr)$. Then
\[
  \Uf\ket{x}\ket{-}
    =\ket{x}\,\tfrac{1}{\sqrt2}\bigl(\ket{0\oplus f(x)}-\ket{1\oplus f(x)}\bigr)
    =(-1)^{f(x)}\ket{x}\ket{-}.
\]
The answer qubit is unchanged and factors out; on the $n$-qubit register the
oracle acts as $\ket{x}\mapsto(-1)^{f(x)}\ket{x}$. This phase-kickback trick lets
us treat the oracle as an $n$-qubit operator that simply attaches a minus sign to
the marked basis state.

\section{The Two Reflections and the Grover Iterate}

Two operators drive the algorithm. Both are \emph{reflections}, and their product
is a rotation. We first define them, prove they are genuine reflections, and then
assemble the Grover iterate.

Throughout, let
\[
  \ket{\psi_0}=H^{\otimes n}\ket{0^{n}}=\frac{1}{\sqrt N}\sum_{x\in\{0,1\}^{n}}\ket{x}
\]
be the uniform superposition over all $N$ basis states. This is the algorithm's
starting state; every candidate carries amplitude $1/\sqrt N$.

\begin{definition}[Marking reflection and diffusion operator]
The \emph{marking reflection} (phase oracle) is
\[
  R_{x_0}=\id-2\ketbra{x_0}{x_0},
\]
which flips the sign of the amplitude on $\ket{x_0}$ and leaves the orthogonal
complement fixed. The \emph{diffusion operator} is
\[
  R_{\psi_0}=2\ketbra{\psi_0}{\psi_0}-\id,
\]
which fixes $\ket{\psi_0}$ and flips the sign of everything orthogonal to it.
It can be written equivalently as
\[
  R_{\psi_0}=H^{\otimes n}\bigl(2\ketbra{0^{n}}{0^{n}}-\id\bigr)H^{\otimes n},
\]
which is how it is implemented as a circuit.
\end{definition}

\begin{lemma}[Both operators are reflections]
\label{l4:lem:reflections}
Each of $R_{x_0}$ and $R_{\psi_0}$ is Hermitian, unitary, and an involution
($R^2=\id$); geometrically each is a reflection of the Hilbert space about a
hyperplane.
\end{lemma}

\begin{proof}
Let $\ket{u}$ be a unit vector and set $P=\ketbra{u}{u}$, so $P=P^{\dagger}$ and
$P^{2}=P$ (an orthogonal projector). Consider $R=\id-2P$. Then
\[
  R^{\dagger}=\id-2P^{\dagger}=\id-2P=R,
\]
so $R$ is Hermitian. Moreover
\[
  R^{2}=(\id-2P)^{2}=\id-4P+4P^{2}=\id-4P+4P=\id,
\]
so $R$ is an involution. A Hermitian involution is unitary, since
$R^{\dagger}R=R^{2}=\id$. Its action is $R\ket{u}=-\ket{u}$ and $R\ket{w}=\ket{w}$
for every $\ket{w}\perp\ket{u}$, which is exactly a reflection about the hyperplane
$\ket{u}^{\perp}$.

Applying this with $\ket{u}=\ket{x_0}$ gives $R_{x_0}=\id-2\ketbra{x_0}{x_0}$
directly. For the diffusion operator, note
$R_{\psi_0}=2\ketbra{\psi_0}{\psi_0}-\id=-\bigl(\id-2\ketbra{\psi_0}{\psi_0}\bigr)$;
by the computation above $\id-2\ketbra{\psi_0}{\psi_0}$ is a Hermitian unitary
involution, and negating it preserves all three properties. Explicitly
$R_{\psi_0}\ket{\psi_0}=\ket{\psi_0}$ and $R_{\psi_0}\ket{w}=-\ket{w}$ for
$\ket{w}\perp\ket{\psi_0}$, so $R_{\psi_0}$ is the reflection through the axis
$\ket{\psi_0}$.
The identity
$R_{\psi_0}=H^{\otimes n}(2\ketbra{0^{n}}{0^{n}}-\id)H^{\otimes n}$ follows because
$H^{\otimes n}$ is unitary and $H^{\otimes n}\ket{0^{n}}=\ket{\psi_0}$.
\end{proof}

Intuitively $R_{\psi_0}$ implements an ``inversion about the mean'': writing a
state as $\sum_x a_x\ket{x}$ with mean amplitude $\bar a=\tfrac1N\sum_x a_x$, one
checks that $R_{\psi_0}\sum_x a_x\ket{x}=\sum_x(2\bar a-a_x)\ket{x}$, reflecting
each amplitude about $\bar a$.

\begin{definition}[Grover iterate]
The \emph{Grover iterate} is the composition of the two reflections,
\[
  G=R_{\psi_0}R_{x_0}
   =\bigl(2\ketbra{\psi_0}{\psi_0}-\id\bigr)\bigl(\id-2\ketbra{x_0}{x_0}\bigr).
\]
One iteration first marks the solution (via $R_{x_0}$) and then reflects about the
mean (via $R_{\psi_0}$).
\end{definition}

\section{The Invariant Plane and the Rotation Picture}

The power of the algorithm comes from the fact that $G$ never leaves a fixed
two-dimensional plane, inside which it acts as an ordinary rotation.

Define the normalized state orthogonal to $\ket{x_0}$ that lives in the ``bad''
subspace:
\[
  \ket{\psi_0^{\perp}}
    =\frac{1}{\sqrt{N-1}}\sum_{x\neq x_0}\ket{x}.
\]
Then $\braket{x_0|\psi_0^{\perp}}=0$, and the uniform state decomposes as
\begin{equation}
  \ket{\psi_0}
    =\frac{1}{\sqrt N}\ket{x_0}+\frac{\sqrt{N-1}}{\sqrt N}\ket{\psi_0^{\perp}}
    =\sin\frac{\theta}{2}\,\ket{x_0}+\cos\frac{\theta}{2}\,\ket{\psi_0^{\perp}},
  \label{l4:eq:initdecomp}
\end{equation}
where the angle $\theta\in(0,\pi)$ is defined by
\begin{equation}
  \sin\frac{\theta}{2}=\frac{1}{\sqrt N},
  \qquad
  \theta=2\arcsin\frac{1}{\sqrt N}\approx\frac{2}{\sqrt N}\quad(N\gg1).
  \label{l4:eq:theta}
\end{equation}
So $\ket{\psi_0}$ makes a small angle $\theta/2$ with the axis
$\ket{\psi_0^{\perp}}$.

\begin{theorem}[$G$ is a rotation in the plane $\mathcal{P}$]
\label{l4:thm:rotation}
Let $\mathcal{P}=\operatorname{span}\{\ket{x_0},\ket{\psi_0^{\perp}}\}$. Then $G$
maps $\mathcal{P}$ into itself, and restricted to $\mathcal{P}$ (in the ordered
orthonormal basis $\{\ket{\psi_0^{\perp}},\ket{x_0}\}$) it is the rotation by the
angle $\theta$ of \eqref{l4:eq:theta}:
\[
  G\big|_{\mathcal{P}}
    =\begin{pmatrix}\cos\theta & -\sin\theta\\[2pt]\sin\theta & \cos\theta\end{pmatrix}.
\]
\end{theorem}

\begin{proof}
First note that both reflections preserve $\mathcal{P}$. The vector $\ket{x_0}$
lies in $\mathcal{P}$ and $\ket{\psi_0}\in\mathcal{P}$ by \eqref{l4:eq:initdecomp},
so $\ket{\psi_0^{\perp}}\in\mathcal{P}$ as well. The operator $R_{x_0}$ negates
$\ket{x_0}$ and fixes every vector orthogonal to $\ket{x_0}$; in particular it
fixes $\ket{\psi_0^{\perp}}$, so it maps $\mathcal{P}$ to itself. The operator
$R_{\psi_0}$ fixes $\ket{\psi_0}\in\mathcal P$ and negates its orthogonal
complement; the in-plane vector orthogonal to $\ket{\psi_0}$ inside $\mathcal P$
stays in $\mathcal P$, so $R_{\psi_0}$ also preserves $\mathcal{P}$. Hence
$G=R_{\psi_0}R_{x_0}$ preserves $\mathcal{P}$.

We now use the plane-geometry fact that \emph{the product of two reflections is a
rotation by twice the angle between their axes}. Inside $\mathcal P$, the mirror
line of $R_{x_0}$ is the axis $\ket{\psi_0^{\perp}}$ (the fixed direction of
$R_{x_0}$ within $\mathcal P$), and the mirror line of $R_{\psi_0}$ is the axis
$\ket{\psi_0}$. By \eqref{l4:eq:initdecomp} the angle between $\ket{\psi_0^{\perp}}$
and $\ket{\psi_0}$ is $\theta/2$, so $G$ is a rotation by $2\cdot(\theta/2)=\theta$.

For a self-contained verification, compute the matrix directly. Using
$\sin\tfrac{\theta}{2}=1/\sqrt N$, $\cos\tfrac{\theta}{2}=\sqrt{(N-1)/N}$ and
\eqref{l4:eq:initdecomp}, apply $R_{x_0}=\id-2\ketbra{x_0}{x_0}$:
\[
  R_{x_0}\ket{\psi_0^{\perp}}=\ket{\psi_0^{\perp}},
  \qquad
  R_{x_0}\ket{x_0}=-\ket{x_0}.
\]
Thus in the basis $\{\ket{\psi_0^{\perp}},\ket{x_0}\}$,
$R_{x_0}=\begin{psmallmatrix}1&0\\0&-1\end{psmallmatrix}$, a reflection about the
$\ket{\psi_0^{\perp}}$ axis. Next $R_{\psi_0}=2\ketbra{\psi_0}{\psi_0}-\id$ is the
reflection about the axis $\ket{\psi_0}$, which sits at angle $\theta/2$ from
$\ket{\psi_0^{\perp}}$; a reflection about a line at angle $\phi$ from the first
basis vector has matrix
$\begin{psmallmatrix}\cos2\phi&\sin2\phi\\\sin2\phi&-\cos2\phi\end{psmallmatrix}$,
so with $\phi=\theta/2$,
\[
  R_{\psi_0}=\begin{pmatrix}\cos\theta&\sin\theta\\ \sin\theta&-\cos\theta\end{pmatrix}.
\]
Multiplying,
\[
  G=R_{\psi_0}R_{x_0}
   =\begin{pmatrix}\cos\theta&\sin\theta\\ \sin\theta&-\cos\theta\end{pmatrix}
    \begin{pmatrix}1&0\\0&-1\end{pmatrix}
   =\begin{pmatrix}\cos\theta&-\sin\theta\\ \sin\theta&\cos\theta\end{pmatrix},
\]
which is exactly the rotation by $\theta$.
\end{proof}

\begin{theorem}[State after $k$ iterations]
\label{l4:thm:afterk}
With $\theta$ as in \eqref{l4:eq:theta},
\[
  G^{k}\ket{\psi_0}
    =\sin\!\frac{(2k+1)\theta}{2}\,\ket{x_0}
     +\cos\!\frac{(2k+1)\theta}{2}\,\ket{\psi_0^{\perp}}.
\]
Consequently the probability of measuring the marked item after $k$ iterations is
$p_k=\sin^{2}\!\bigl((2k+1)\theta/2\bigr)$.
\end{theorem}

\begin{proof}
By Theorem~\ref{l4:thm:rotation}, $G$ restricted to $\mathcal P$ is rotation by
$\theta$, so $G^{k}$ is rotation by $k\theta$. The initial state
\eqref{l4:eq:initdecomp} sits at angle $\theta/2$ measured from the
$\ket{\psi_0^{\perp}}$ axis (its $\ket{x_0}$-component is
$\sin(\theta/2)$). Rotating by $k\theta$ increases this angle to
$\theta/2+k\theta=(2k+1)\theta/2$. Hence the $\ket{x_0}$-component becomes
$\sin\bigl((2k+1)\theta/2\bigr)$ and the $\ket{\psi_0^{\perp}}$-component becomes
$\cos\bigl((2k+1)\theta/2\bigr)$, as claimed. Measuring in the computational basis
yields $x_0$ with probability equal to the squared $\ket{x_0}$-amplitude, namely
$p_k=\sin^{2}\bigl((2k+1)\theta/2\bigr)$.
\end{proof}

\section{Optimal Iteration Count and Query Complexity}

Theorem~\ref{l4:thm:afterk} shows the amplitude on $\ket{x_0}$ is maximized when the
angle $(2k+1)\theta/2$ is as close as possible to $\pi/2$, i.e.\ when the state is
(nearly) aligned with $\ket{x_0}$.

\begin{theorem}[Optimal number of iterations]
\label{l4:thm:kstar}
Choose $k$ so that $(2k+1)\theta/2$ is closest to $\pi/2$; the optimal choice is
\[
  k^{\ast}=\left\lfloor\frac{\pi}{4}\sqrt N\right\rfloor
  \qquad\bigl(\text{rounding } \tfrac{\pi}{2\theta}-\tfrac12\bigr).
\]
With this choice the success probability satisfies $p_{k^{\ast}}\ge 1-1/N$, so a
single run succeeds with probability $\approx1$. The algorithm therefore uses
$k^{\ast}=O(\sqrt N)$ oracle queries.
\end{theorem}

\begin{proof}
We want $(2k+1)\theta/2\approx\pi/2$, i.e.\ $2k+1\approx\pi/\theta$, i.e.\
$k\approx\pi/(2\theta)-1/2$. Using the approximation
$\theta\approx2/\sqrt N$ from \eqref{l4:eq:theta} valid for large $N$,
\[
  k\approx\frac{\pi}{2\theta}-\frac12\approx\frac{\pi}{2}\cdot\frac{\sqrt N}{2}-\frac12
     =\frac{\pi}{4}\sqrt N-\frac12,
\]
and rounding to the nearest nonnegative integer gives
$k^{\ast}=\lfloor(\pi/4)\sqrt N\rfloor$. Because $k$ is an integer we cannot in
general hit the angle $\pi/2$ exactly, but the residual angular error is at most
$\theta/2$ (half a rotation step): choosing $k^{\ast}$ so that
$(2k^{\ast}+1)\theta/2$ lies within $\theta/2$ of $\pi/2$ gives
\[
  p_{k^{\ast}}=\sin^{2}\!\frac{(2k^{\ast}+1)\theta}{2}
    =\cos^{2}\!\Bigl(\frac{\pi}{2}-\frac{(2k^{\ast}+1)\theta}{2}\Bigr)
    \ge\cos^{2}\!\frac{\theta}{2}=1-\sin^{2}\frac{\theta}{2}=1-\frac1N.
\]
Each iterate $G=R_{\psi_0}R_{x_0}$ makes exactly one oracle call (inside
$R_{x_0}$), so the total query cost is $k^{\ast}=O(\sqrt N)$.
\end{proof}

This is the headline result: $O(\sqrt N)$ queries versus the classical $O(N)$, a
quadratic speedup. The gain comes from operating on \emph{amplitudes} rather than
probabilities: $G$ rotates the amplitude linearly in $k$, and probability is its
square, so the linear growth of the angle produces the square-root scaling in $k$.

\begin{remark}[Do not over-rotate]
More iterations is \emph{not} always better. Since $p_k=\sin^{2}((2k+1)\theta/2)$
oscillates, once $(2k+1)\theta/2$ passes $\pi/2$ the success probability decreases
again --- the state ``rotates past'' $\ket{x_0}$. One must therefore stop at
$k^{\ast}$; running the loop too long is counterproductive. If $N$ is unknown one
can still succeed with $O(\sqrt N)$ expected queries by trying geometrically
increasing guesses for the number of iterations and choosing the number at random
within each guess.
\end{remark}

\section{Multiple Solutions and Amplitude Amplification}

Suppose now that there are $M\ge1$ marked items, $S=\{x:f(x)=1\}$ with $\abs{S}=M$.
Everything goes through with the ``good'' subspace replaced by the span of the
marked basis states.

Let
\[
  \ket{\mathrm{good}}=\frac{1}{\sqrt M}\sum_{x\in S}\ket{x},
  \qquad
  \ket{\mathrm{bad}}=\frac{1}{\sqrt{N-M}}\sum_{x\notin S}\ket{x},
\]
so that $\ket{\psi_0}=\sqrt{M/N}\,\ket{\mathrm{good}}+\sqrt{(N-M)/N}\,\ket{\mathrm{bad}}$.
Define $R_S=\id-2\sum_{x\in S}\ketbra{x}{x}$ (the phase oracle now flips all marked
items) and keep the same diffusion $R_{\psi_0}$. The identical argument shows
$G=R_{\psi_0}R_S$ is a rotation by $\theta$ in
$\operatorname{span}\{\ket{\mathrm{good}},\ket{\mathrm{bad}}\}$ with
\[
  \sin\frac{\theta}{2}=\sqrt{\frac{M}{N}},
  \qquad\text{equivalently}\qquad
  \sin\theta=2\sin\frac{\theta}{2}\cos\frac{\theta}{2}
            =2\sqrt{\frac{M}{N}}\sqrt{\frac{N-M}{N}}\approx 2\sqrt{\frac{M}{N}}.
\]
The optimal iteration count becomes
\[
  k^{\ast}=\left\lfloor\frac{\pi}{4}\sqrt{\frac{N}{M}}\right\rfloor
  =O\!\left(\sqrt{\frac{N}{M}}\right),
\]
so more solutions make the search faster.

\begin{remark}[Amplitude amplification]
Grover's construction generalizes far beyond search. Suppose $\mathcal A$ is any
(quantum) procedure preparing a state
$\mathcal A\ket{0}=\sqrt{p_0}\,\ket{\psi_{\mathrm{good}}}+\sqrt{1-p_0}\,\ket{\psi_{\mathrm{bad}}}$,
where a projector $P$ recognizes the good part and $p_0$ is the initial success
probability. Replacing $R_{x_0}$ by the reflection $\id-2P$ and $R_{\psi_0}$ by
$\mathcal A(2\ketbra{0}{0}-\id)\mathcal A^{\dagger}$ yields an iterate that rotates
by an angle $\theta$ with $\sin^{2}(\theta/2)=p_0$. After $O(1/\sqrt{p_0})$
iterations the good state is reached with probability $\Omega(1)$ --- a quadratic
improvement over the $O(1/p_0)$ repetitions that classical
prepare-and-check would require. Grover search is the special case
$p_0=M/N$.
\end{remark}

\begin{example}[The exact case $N=4$, $M=1$]
Here $\sin(\theta/2)=1/\sqrt4=1/2$, so $\theta/2=\pi/6$ and $\theta=\pi/3$. After a
single iteration ($k=1$) the angle is
\[
  \frac{(2\cdot1+1)\theta}{2}=\frac{3\theta}{2}=\frac{3}{2}\cdot\frac{\pi}{3}=\frac{\pi}{2},
\]
so by Theorem~\ref{l4:thm:afterk} the state is
$G\ket{\psi_0}=\sin(\pi/2)\ket{x_0}+\cos(\pi/2)\ket{\psi_0^{\perp}}=\ket{x_0}$.
Thus for $N=4$ a \emph{single} Grover iteration finds the marked item with
certainty --- a rare instance where the algorithm is exactly deterministic.
\end{example}

\begin{example}[A multi-solution estimate]
Suppose a SAT instance on $n=30$ variables has $M=2^{10}$ satisfying assignments
among $N=2^{30}$ candidates. Classical exhaustive search needs on the order of
$N/M=2^{20}\approx10^{6}$ trials to hit a solution. Grover instead needs about
\[
  \frac{\pi}{4}\sqrt{\frac{N}{M}}
    =\frac{\pi}{4}\sqrt{2^{20}}
    =\frac{\pi}{4}\cdot2^{10}\approx804
\]
oracle queries --- the square-root improvement in action.
\end{example}

\section{Optimality: the \texorpdfstring{$\Omega(\sqrt N)$}{Omega(sqrt N)} Lower Bound}

Grover's algorithm is not just good; it is asymptotically the best possible.

\begin{theorem}[Query lower bound]
Any quantum algorithm that finds the marked item $x_0$ (or even merely decides
whether a marked item exists) with constant success probability must make
$\Omega(\sqrt N)$ oracle queries.
\end{theorem}

\begin{proof}[Idea]
Run the algorithm both on the ``empty'' oracle $f\equiv0$ and on the oracle marking
a single item $x_0$. Let $\ket{\phi_k}$ be the state after $k$ queries in the
empty case and $\ket{\phi_k^{x_0}}$ the state in the marked case, and define the
total deviation
\[
  D_k=\sum_{x_0}\norm{\,\ket{\phi_k^{x_0}}-\ket{\phi_k}\,}^{2}.
\]
A hybrid argument shows each additional query can increase $D_k$ by only a bounded
amount: $D_k\le 4k^{2}$. On the other hand, to reliably distinguish which item is
marked, the marked-case states for different $x_0$ must be nearly orthogonal to the
empty-case state, which forces $D_k=\Omega(N)$. Combining
$\Omega(N)\le D_k\le 4k^{2}$ gives $k=\Omega(\sqrt N)$.
\end{proof}

\begin{remark}[Grover is optimal]
The upper bound of Theorem~\ref{l4:thm:kstar} and the lower bound above match up to
constant factors: both are $\Theta(\sqrt N)$. Hence Grover's algorithm is optimal
for unstructured search --- no quantum algorithm can do asymptotically better. The
quadratic speedup is the best achievable when the only access to the problem is
through an unstructured oracle; a super-polynomial speedup would require exploitable
structure, as in Shor's algorithm or the phase-estimation-based methods.
\end{remark}

\section{Problems}

\begin{exercise}
\label{l4:ex:invmean}
Show that the diffusion operator implements ``inversion about the mean'': for any
state $\sum_x a_x\ket{x}$ with mean amplitude $\bar a=\tfrac1N\sum_x a_x$,
\[
  R_{\psi_0}\sum_x a_x\ket{x}=\sum_x(2\bar a-a_x)\ket{x}.
\]
\end{exercise}

\begin{exercise}
\label{l4:ex:phaseoracle}
Verify the phase-kickback identity
\[
  \Uf\ket{x}\ket{-}=(-1)^{f(x)}\ket{x}\ket{-}
\]
starting from $\Uf\ket{x}\ket{y}=\ket{x}\ket{y\oplus f(x)}$, and explain why this
lets us regard the oracle as the $n$-qubit operator $R_{x_0}$.
\end{exercise}

\begin{exercise}
\label{l4:ex:n2}
Take $N=2$ (one qubit, $n=1$) with the single marked item $x_0$. Compute $\theta$
from \eqref{l4:eq:theta}, apply one Grover iterate to $\ket{\psi_0}$, and interpret the
result. Why is the case $N=2$ slightly special?
\end{exercise}

\begin{exercise}
\label{l4:ex:success}
For general $N$, using $k^{\ast}=\lfloor(\pi/4)\sqrt N\rfloor$, show that the
success probability satisfies $p_{k^{\ast}}\ge1-1/N$. (You may use that the
residual angle from $\pi/2$ is at most $\theta/2$.)
\end{exercise}

\begin{exercise}
\label{l4:ex:overrotate}
Consider $N=2^{20}$. Approximately how many iterations $k^{\ast}$ maximize the
success probability? What happens qualitatively to the success probability if one
instead runs $2k^{\ast}$ iterations? Explain using $p_k=\sin^2((2k+1)\theta/2)$.
\end{exercise}

\begin{exercise}
\label{l4:ex:multi}
There are $M=4$ marked items among $N=1024$. Find $\sin(\theta/2)$, the optimal
iteration count $k^{\ast}$, and compare with the single-solution count for the same
$N$.
\end{exercise}

\begin{exercise}
\label{l4:ex:reflprod}
Prove the plane-geometry fact used in Theorem~\ref{l4:thm:rotation}: in a real
two-dimensional plane, the composition of a reflection about a line $\ell_1$ and a
reflection about a line $\ell_2$ (with $\ell_2$ at angle $\alpha$ from $\ell_1$) is
a rotation by $2\alpha$.
\end{exercise}

\subsection*{Solutions}

\begin{solution}[Exercise~\ref{l4:ex:invmean}]
Write $R_{\psi_0}=2\ketbra{\psi_0}{\psi_0}-\id$ and let $\ket{\phi}=\sum_x a_x\ket{x}$.
Since $\ket{\psi_0}=\tfrac1{\sqrt N}\sum_x\ket{x}$, we have
$\braket{\psi_0|\phi}=\tfrac1{\sqrt N}\sum_x a_x=\sqrt N\,\bar a$. Therefore
\[
  2\ketbra{\psi_0}{\psi_0}\phi
    =2\sqrt N\,\bar a\,\ket{\psi_0}
    =2\sqrt N\,\bar a\cdot\frac1{\sqrt N}\sum_x\ket{x}
    =\sum_x 2\bar a\,\ket{x}.
\]
Subtracting $\id\ket{\phi}=\sum_x a_x\ket{x}$ gives
$R_{\psi_0}\ket{\phi}=\sum_x(2\bar a-a_x)\ket{x}$, i.e.\ each amplitude is reflected
about the mean $\bar a$. This is why $R_{\psi_0}$ is often called ``inversion about
the mean.''
\end{solution}

\begin{solution}[Exercise~\ref{l4:ex:phaseoracle}]
With $\ket{-}=\tfrac1{\sqrt2}(\ket0-\ket1)$,
\[
  \Uf\ket{x}\ket{-}
   =\tfrac1{\sqrt2}\bigl(\ket{x}\ket{0\oplus f(x)}-\ket{x}\ket{1\oplus f(x)}\bigr).
\]
If $f(x)=0$ this is $\tfrac1{\sqrt2}(\ket{x}\ket0-\ket{x}\ket1)=\ket{x}\ket{-}$. If
$f(x)=1$ the two kets swap, giving
$\tfrac1{\sqrt2}(\ket{x}\ket1-\ket{x}\ket0)=-\ket{x}\ket{-}$. Both cases combine into
$(-1)^{f(x)}\ket{x}\ket{-}$. The answer qubit stays in $\ket{-}$ and factors out, so
on the $n$-qubit register the oracle acts as $\ket{x}\mapsto(-1)^{f(x)}\ket{x}$,
which negates exactly $\ket{x_0}$ --- that is $R_{x_0}=\id-2\ketbra{x_0}{x_0}$.
\end{solution}

\begin{solution}[Exercise~\ref{l4:ex:n2}]
For $N=2$, $\sin(\theta/2)=1/\sqrt2$, so $\theta/2=\pi/4$ and $\theta=\pi/2$. The
initial state is $\ket{\psi_0}=\tfrac1{\sqrt2}(\ket{x_0}+\ket{\psi_0^{\perp}})$, at
angle $\theta/2=\pi/4$ from the $\ket{\psi_0^{\perp}}$ axis. One iterate rotates by
$\theta=\pi/2$, bringing the angle to $\pi/4+\pi/2=3\pi/4$; equivalently
$(2\cdot1+1)\theta/2=3\pi/4$, and $p_1=\sin^2(3\pi/4)=1/2$. So a single iterate does
\emph{not} give certainty here, and indeed
$\lfloor(\pi/4)\sqrt2\rfloor=\lfloor1.11\rfloor=1$ while the best achievable
probability for $N=2$ is $1/2$, no better than a random guess. The case is special
because with only two items the initial angle $\pi/4$ is already large and the
rotation overshoots; unstructured search offers no advantage when $N=2$.
\end{solution}

\begin{solution}[Exercise~\ref{l4:ex:success}]
Pick $k^{\ast}$ so that the angle $(2k^{\ast}+1)\theta/2$ is within $\theta/2$ of
$\pi/2$; this is possible because consecutive angles differ by $\theta$. Writing the
residual $\delta=\pi/2-(2k^{\ast}+1)\theta/2$ with $\abs{\delta}\le\theta/2$,
\[
  p_{k^{\ast}}=\sin^2\!\Bigl(\tfrac\pi2-\delta\Bigr)=\cos^2\delta\ge\cos^2\tfrac\theta2
   =1-\sin^2\tfrac\theta2=1-\tfrac1N,
\]
using $\sin(\theta/2)=1/\sqrt N$. Hence the success probability is at least
$1-1/N\to1$, so one run essentially always succeeds; if it fails one verifies with a
single extra oracle call and repeats.
\end{solution}

\begin{solution}[Exercise~\ref{l4:ex:overrotate}]
For $N=2^{20}=1{,}048{,}576$, $\sqrt N=1024$, so
$k^{\ast}=\lfloor(\pi/4)\cdot1024\rfloor=\lfloor 804.2\rfloor=804$. At $k^{\ast}$ the
angle $(2k^{\ast}+1)\theta/2\approx\pi/2$ and $p_{k^{\ast}}\approx1$. Running instead
$2k^{\ast}\approx1608$ iterations roughly doubles the angle to about
$(2\cdot2k^{\ast}+1)\theta/2\approx 3\pi/2$; since $p_k$ oscillates as $\sin^2$, the
state has rotated well past $\ket{x_0}$ and the success probability has fallen back
toward small values. The lesson: the probability oscillates in $k$, so over-rotating
destroys the amplitude gained --- stop at $k^{\ast}$.
\end{solution}

\begin{solution}[Exercise~\ref{l4:ex:multi}]
With $M=4$, $N=1024$: $\sin(\theta/2)=\sqrt{M/N}=\sqrt{4/1024}=\sqrt{1/256}=1/16$.
The optimal count is
$k^{\ast}=\lfloor(\pi/4)\sqrt{N/M}\rfloor=\lfloor(\pi/4)\sqrt{256}\rfloor
=\lfloor(\pi/4)\cdot16\rfloor=\lfloor12.57\rfloor=12$. For a single solution
($M=1$) with the same $N$ one would need
$\lfloor(\pi/4)\sqrt{1024}\rfloor=\lfloor(\pi/4)\cdot32\rfloor=\lfloor25.13\rfloor=25$
iterations. So having $M=4$ solutions roughly halves the query count (from $25$ to
$12$), consistent with the $O(\sqrt{N/M})$ scaling: a factor of $2$ speedup for a
factor of $4$ more solutions.
\end{solution}

\begin{solution}[Exercise~\ref{l4:ex:reflprod}]
Place coordinates so that $\ell_1$ is the $x$-axis and $\ell_2$ makes angle $\alpha$
with it. Reflection about a line at angle $\beta$ has matrix
$M(\beta)=\begin{psmallmatrix}\cos2\beta&\sin2\beta\\\sin2\beta&-\cos2\beta\end{psmallmatrix}$.
Thus reflection about $\ell_1$ is
$M(0)=\begin{psmallmatrix}1&0\\0&-1\end{psmallmatrix}$
and about $\ell_2$ is $M(\alpha)$. Their product is
\[
  M(\alpha)M(0)
   =\begin{pmatrix}\cos2\alpha&\sin2\alpha\\ \sin2\alpha&-\cos2\alpha\end{pmatrix}
    \begin{pmatrix}1&0\\0&-1\end{pmatrix}
   =\begin{pmatrix}\cos2\alpha&-\sin2\alpha\\ \sin2\alpha&\cos2\alpha\end{pmatrix},
\]
which is the rotation matrix by angle $2\alpha$. Hence composing the two reflections
gives a rotation by twice the angle between the mirror lines. In
Theorem~\ref{l4:thm:rotation} the two mirror axes are $\ket{\psi_0^{\perp}}$ and
$\ket{\psi_0}$, separated by $\theta/2$, giving rotation by $\theta$.
\end{solution}
